% !TEX root = ../main.tex
\chapter{KESIMPULAN DAN SARAN}
\label{chap:kesimpulan}

% ============================================================
\section{Kesimpulan}
\label{sec:kesimpulan}
% ============================================================

Penelitian ini merancang, mengimplementasikan, dan mengevaluasi arsitektur sistem cerdas berbasis komputasi tepi (\textit{edge computing}) dan analitik data besar (\textit{big data analytics}) untuk pelacakan pelanggaran penggunaan helm di Institut Teknologi Sumatera (Itera). Evaluasi kinerja sistem secara menyeluruh menghasilkan tiga kesimpulan utama yang menjawab tujuan penelitian.

\begin{enumerate}[1.]
  \item  Model deteksi objek berbasis YOLOv8 dioptimalkan dengan dataset kustom CCTV kampus dan mencapai \textit{Mean Average Precision} (mAP@50) sebesar 91,4\% pada data validasi serta 90,0\% pada data uji. Selisih marjinal antara \textit{precision} (87,5\%) dan \textit{recall} (86,5\%) memperlihatkan model yang seimbang antara sensitivitas lokalisasi dan presisi pemisahan kelas \texttt{MOTORCYCLE}, \texttt{DRIVER\_HELMET}, dan \texttt{DRIVER\_NO\_HELMET} di tengah variasi distorsi visual dan pencahayaan luar ruang.

  \item Integrasi inferensi YOLOv8 dan pelacakan multiobjek ByteTrack mengubah deteksi piksel mentah menjadi aliran peristiwa (\textit{event stream}) terstruktur. Filter asosiasi spasial (\textit{Intersection over Area}) dan konfirmasi temporal \textit{N-Frame} menggugurkan deteksi positif palsu (\textit{false positives}) sebelum data ditransmisikan. Pada pengujian operasional 10 jam, rangkaian filter tersebut menyaring 11.172 deteksi mentah menjadi 113 \textit{event} pelanggaran yang dipersistenkan ke \textit{data mart}. Arsitektur pemrosesan tepi mencatat \textit{throughput} stabil 18,9 FPS dan latensi inferensi 53 milidetik tanpa dukungan GPU. Distribusi data melalui Apache Kafka (latensi transmisi median 2,65 milidetik) dan Apache Spark Structured Streaming berjalan asinkron tanpa memicu \textit{backpressure} antrian, sehingga sistem tetap stabil pada lonjakan aliran data.

  \item Perancangan \textit{data mart} berbasis skema bintang (\textit{star schema}) pada PostgreSQL menyediakan landasan komputasi untuk \textit{Online Analytical Processing} (OLAP). Pengujian operasional pada data riil selama 10 jam (07.00--17.00 WIB) mengekstraksi wawasan dari 113 kejadian pelanggaran aktual. \textit{Dashboard} Grafana memetakan distribusi kepadatan bimodal, dengan anomali puncak absolut tercatat pada pukul 14.45--15.15 WIB yang berkorelasi dengan faktor cuaca (pasca-hujan deras) dan transisi jadwal kelas. Implementasi infrastruktur ujung-ke-ujung (\textit{end-to-end}) tersebut menempatkan sistem bukan sekadar sebagai alat pendeteksi, melainkan sebagai instrumen Sistem Pendukung Keputusan (\textit{Decision Support System}) yang menggeser pengawasan UPT K3L Itera dari inspeksi reaktif menuju penjagaan preventif berbasis bukti (\textit{data-driven}).
\end{enumerate}

% ============================================================
\section{Saran}
\label{sec:saran}
% ============================================================

Sistem sudah beroperasi sesuai batasan penelitian, namun masih terbuka peluang optimasi agar arsitektur ini siap diterapkan pada skala produksi (\textit{production-grade}). Rekomendasi untuk penelitian lanjutan disusun ke dalam tujuh arah pengembangan berikut.

\begin{enumerate}[1.]
  \item \textbf{Validasi \textit{end-to-end} berbasis pencocokan manual.}
    Evaluasi kinerja sistem saat ini bertumpu pada metrik per komponen (mAP, \textit{throughput}, latensi). Pengembangan ke depan perlu menambahkan validasi \textit{end-to-end} dengan membandingkan \textit{event} keluaran \textit{pipeline} terhadap anotasi manual oleh pengamat manusia pada segmen \textit{peak hour}. Dari pencocokan satu-ke-satu tersebut dapat diturunkan \textit{precision}, \textit{recall}, F1-Score, \textit{counting accuracy}, dan MAPE per jendela 15 menit yang menggambarkan ketepatan sistem sebagai instrumen pencacah pelanggaran, bukan sekadar pendeteksi objek. Praktik pencocokan ini lazim digunakan pada penelitian \textit{video analytics} lapangan dan memperkuat klaim kelayakan operasional.

  \item \textbf{Akselerasi perangkat keras pada simpul tepi (\textit{Edge AI}).}
    Beban inferensi spasial saat ini masih bertumpu pada Unit Pemroses Sentral (CPU). Pengembangan ke arsitektur multikamera (\textit{multi-camera surveillance}) beresolusi tinggi sebaiknya memindahkan inferensi ke akselerator khusus seperti NVIDIA Jetson atau Google Coral \textit{Tensor Processing Unit} (TPU). Akselerator tersebut berpotensi mendorong \textit{throughput} melampaui 30 FPS sekaligus menurunkan latensi operasional secara berarti.

  \item \textbf{Perluasan ekstraksi atribut keselamatan dan identifikasi entitas.}
    Model visi komputer dapat dikembangkan untuk menjangkau pelanggaran sekunder yang juga berkontribusi pada risiko fatalitas, misalnya kelebihan kapasitas penumpang (berbonceng tiga) atau penggunaan ponsel saat berkendara. Integrasi dengan subsistem \textit{Automatic License Plate Recognition} (ALPR) juga relevan agar setiap \textit{event} terikat langsung pada identitas registrasi kendaraan pelanggar dan dapat ditindaklanjuti oleh unit pengamanan kampus.

  \item \textbf{Analitik prediktif dan deteksi dini lonjakan pelanggaran.}
    Catatan operasional 10 jam memperlihatkan pola bimodal pagi--sore yang konsisten dengan jadwal kuliah. Lapisan analitik dapat diperkaya dengan model deret waktu sederhana (mis. SARIMA atau \textit{exponential smoothing}) untuk memproyeksikan lonjakan beban pada jam berikutnya, sehingga UPT K3L menerima peringatan dini sebelum periode rawan dimulai dan dapat menggeser penempatan personel patroli secara proaktif.

  \item \textbf{Penerapan multi-kamera dan fusi sudut pandang.}
    Penelitian ini diuji pada satu kamera gerbang utama. Replikasi pada beberapa titik kritis (gerbang sekunder, persimpangan internal, jalur sepeda motor) bersama mekanisme \textit{re-identification} antar kamera berpeluang memetakan jejak pengendara lintas titik dan mengubah \textit{Track ID} sesi menjadi identitas pelanggar fisik yang persisten---salah satu keterbatasan yang ditandai pada Bab~\ref{chap:hasil}.

  \item \textbf{Penerapan praktik MLOps untuk siklus hidup model.}
    Pelatihan model pada penelitian ini berlangsung satu kali di platform Roboflow Train 3.0, sedangkan operasi produksi memerlukan pembaruan berkelanjutan agar akurasi tidak menurun seiring perubahan musim, sudut kamera, atau ragam kendaraan baru. Penelitian lanjutan dapat membangun jalur MLOps yang menggabungkan registri model (\textit{model registry}), pelacakan versi dataset, dan orkestrasi pelatihan otomatis melalui kakas seperti MLflow, DVC, dan Apache Airflow. Pemantauan pergeseran distribusi (\textit{model drift}) pada keluaran produksi dapat dijalankan dengan kakas seperti Evidently AI atau WhyLabs sehingga pelatihan ulang dipicu otomatis ketika metrik turun di bawah ambang yang ditetapkan. Mekanisme \textit{human-in-the-loop} pada anotasi \textit{event} berkeyakinan rendah menjaga mutu label sekaligus memperkaya dataset secara berkelanjutan, sehingga arsitektur tumbuh dari purwarupa satu kali pelatihan menjadi sistem yang adaptif terhadap perubahan lapangan.

  \item \textbf{Perluasan jalur deteksi ke jenis pelanggaran lalu lintas lain.}
    Sistem saat ini fokus pada pelanggaran helm, namun arsitektur kelas YOLO yang berpasangan dengan filter spasial-temporal sebenarnya siap menampung jenis pelanggaran lain dengan biaya pengembangan yang relatif rendah. Penelitian lanjutan dapat menambahkan jalur deteksi untuk pelanggaran pergerakan seperti melawan arus, motor masuk jalur pejalan kaki, parkir liar di luar petak, dan penerobosan rambu larangan di persimpangan internal kampus. Pelanggaran pergerakan memerlukan analisis lintasan (\textit{trajectory}) lewat \texttt{Track ID} ByteTrack yang dibandingkan terhadap poligon zona dan arah lajur yang ditetapkan operator, bukan asosiasi IoA seperti pada kasus helm. Pengembangan dataset menjadi titik kritis karena setiap kategori membutuhkan anotasi tersendiri; pemanfaatan dataset publik seperti Indian Driving Dataset (IDD) atau dataset CCTV lalu lintas terbuka dapat menjadi \textit{bootstrap} sebelum dataset kampus terkumpul. Dengan struktur yang sama, satu \textit{pipeline} sisi tepi dapat memantau beberapa jenis pelanggaran sekaligus, dan \textit{dashboard} Grafana hanya perlu menambah panel per kategori tanpa mengubah skema bintang \textit{data mart}.
\end{enumerate}

Secara keseluruhan, penelitian ini menyajikan kerangka kerja arsitektur analitik video terintegrasi yang memadukan pembelajaran mendalam (YOLOv8), pelacakan multiobjek (ByteTrack), dan pemrosesan aliran data (Kafka--Spark) sebagai solusi pemantauan keselamatan lalu lintas cerdas di lingkungan perguruan tinggi. Kombinasi modul deteksi, filter spasial-temporal, dan \textit{data mart} bintang menggeser pengawasan helm dari catatan manual menjadi rangkaian \textit{event} terukur yang siap dijadikan basis kebijakan. Ketujuh arah pengembangan di atas membuka jalur kelanjutan agar arsitektur ini tumbuh dari purwarupa satu kamera menjadi infrastruktur kampus yang mandiri dan adaptif.
